\begin{frame}{Alterations}
	\noindent
	
	\begin{itemize}
		\item Different representation theorem
		\begin{center}
			\renewcommand{\arraystretch}{1.8}
			\setlength{\tabcolsep}{8pt}
			\begin{table}[ht]
				\centering
				\resizebox{0.98\textwidth}{!}{%
					\begin{tabular}{|p{7.5cm}|p{8.0cm}|p{3.3cm}|}
						
						\rowcolor{brown!60}
						\hline
						\textbf{Name} & \textbf{TLDR} & \textbf{Reference} \\ \hline
						Deep Learning Alternatives of the
						Kolmogorov Superposition Theorem & Utilizes the Laczkovich Representation Theorem to reduce the number of learned parameters &  Guilhoto and  Perdikaris (2025)\\ \hline
					\end{tabular}
				}%
			\end{table}
		\end{center}
	\end{itemize}
\end{frame}

\begin{frame}{Alterations}
	\noindent
	\begin{itemize}
		\item Different internal representation
		\begin{center}
			\renewcommand{\arraystretch}{1.8}
			\setlength{\tabcolsep}{8pt}
			\begin{table}[ht]
				\centering
				\resizebox{0.98\textwidth}{!}{%
					\begin{tabular}{|p{7.5cm}|p{8.0cm}|p{3.3cm}|}
						\rowcolor{brown!60}
						\hline
						\textbf{Name} & \textbf{TLDR} & \textbf{Reference} \\ \hline
						Chebyshev Polynomial-Based Kolmogorov-Arnold Networks & Activattion functions are modeled with Chebyshev Polynomials & Sidharth  et al. (2024) \\ \hline
						fKAN: Fractional Kolmogorov-Arnold Networks with trainable Jacobi basis functions & Utilizes fractional-orthogonal Jacobi functions as basis functions &  Aghaei (2025)\\ \hline
						Wavelet Kolmogorov-Arnold Networks (Wav-KAN) & Uses wavelets instead of B-splines; Inputs are complete MNIST images & Bozorgasl and Chen \\ \hline
					\end{tabular}
				}%
			\end{table}
		\end{center}
	\end{itemize}
\end{frame}

\begin{frame}{Extensions}
		\begin{itemize}
		\item Temporal
	\noindent
	\begin{center}
		\renewcommand{\arraystretch}{1.8}
		\setlength{\tabcolsep}{8pt}
		\begin{table}[ht]
			\centering
			\resizebox{0.98\textwidth}{!}{%
				\begin{tabular}{|p{6.5cm}|p{9cm}|p{3.3cm}|}
					
					\rowcolor{brown!60}
					\hline
					\textbf{Name} & \textbf{TLDR} & \textbf{Reference} \\ \hline
					Temporal Kolmogorov-Arnold Networks (TKAN) & Changes the output of the LSTM cell to $\sigma(\text{KAN}(x))$ & Genet and Inzirillo (2024) \\ \hline
					Temporal Kolmogorov-Arnold Transformer (TKAT) & An encoder-decoder architecture that utilizes TKAN to both components & Genet and Inzirillo (2024) \\ \hline
					seqKAN   &  Replaces the hidden and output layers of an RNN with pure KAN functions, i.e., without trainable weights   & Boura and Konstantopoulos (2025) \\ \hline
				\end{tabular}
			}%
		\end{table}
	\end{center}
\end{itemize}
\end{frame}


\begin{frame}{Attempts for Efficiency}
	\noindent
	\begin{center}
		\renewcommand{\arraystretch}{1.8}
		\setlength{\tabcolsep}{8pt}
		\begin{table}[ht]
			\centering
			\resizebox{0.98\textwidth}{!}{%
				\begin{tabular}{|p{6.5cm}|p{9cm}|p{3.3cm}|}
					
					\rowcolor{brown!60}
					\hline
					\textbf{Name} & \textbf{TLDR} & \textbf{Reference} \\ \hline
					FastKAN &  & \\ \hline
					FBKANs &  & \\ \hline
					DropKAN &  & \\ \hline
				\end{tabular}
			}%
		\end{table}
	\end{center}
\end{frame}


\begin{frame}{Extensions}
	\begin{itemize}
		\item Graph
	
	\noindent
	\vspace{-0.6cm}
	\begin{center}
		\renewcommand{\arraystretch}{1.8}
		\setlength{\tabcolsep}{8pt}
		\begin{table}[ht]
			\centering
			\resizebox{0.98\textwidth}{!}{%
				\begin{tabular}{|p{6.5cm}|p{9cm}|p{3.3cm}|}
					\rowcolor{brown!60}
					\hline
					\textbf{Name} & \textbf{TLDR} & \textbf{Reference} \\ \hline
					GKAN: Graph Kolmogorov-Arnold Networks & Performs the aggregation of the graph nodes with KANs & Kiamari et al. (2024) \\ \hline
				\end{tabular}
			}%
		\end{table}
	\end{center}
	\vspace{0.5cm}
	\item Quantum
	\noindent
	\begin{center}
		\renewcommand{\arraystretch}{1.8}
		\setlength{\tabcolsep}{8pt}
		\begin{table}[ht]
			\centering
			\resizebox{0.98\textwidth}{!}{%
				\begin{tabular}{|p{6.5cm}|p{9cm}|p{3.3cm}|}
					\rowcolor{brown!60}
					\hline
					\textbf{Name} & \textbf{TLDR} & \textbf{Reference} \\ \hline
					QKAN: Quantum Kolmogorov-Arnold Networks & Exploits quantum linear algebra tools to apply parameterized activation functions on the edges of the network & Ivashkov et al. (2024) \\ \hline
				\end{tabular}
			}%
		\end{table}
	\end{center}
\end{itemize}
\end{frame}

\begin{frame}{Applications}
	\noindent
	\begin{center}
		\renewcommand{\arraystretch}{1.8}
		\setlength{\tabcolsep}{8pt}
		\begin{table}[ht]
			\centering
			\resizebox{0.98\textwidth}{!}{%
				\begin{tabular}{|p{3.7cm}|p{5.7cm}|p{6cm}|p{3.3cm}|}
					%\hline
					\rowcolor{brown!60}
					\hline
					\textbf{Domain/Application} & \textbf{Name} & \textbf{TLDR} & \textbf{Reference} \\ \hline
					Physics Informed Modeling & Kolmogorov-Arnold Network Ordinary Differential Equations (KAN-ODEs) & Vanilla KAN used as an approximator for a dynamical system in the neural ODE framework & Koenig et al. (2024)\\ \hline
					Human Activity Recognition & iKAN & Incremental Learning paired with KANs; Encoder and froxen KAN-based classifier training & Liu et al. (2024)\\ \hline
				\end{tabular}
			
			}%
		\end{table}
	\end{center}
\end{frame}


\begin{frame}{ECAI Honorable Mentions}
	\noindent
	\begin{center}
		\renewcommand{\arraystretch}{1.8}
		\setlength{\tabcolsep}{8pt}
		\begin{table}[ht]
			\centering
			\resizebox{0.98\textwidth}{!}{%
				\begin{tabular}{|p{6.5cm}|p{9cm}|p{3.3cm}|}
					\rowcolor{brown!60}
					\hline
					\textbf{Name} & \textbf{TLDR} & \textbf{Reference} \\ \hline
					 &  &  \\ \hline
				\end{tabular}
			}%
		\end{table}
	\end{center}
\end{frame}



